\documentclass[10pt,a4paper]{amsart}
\usepackage[margin=19mm]{geometry}
\usepackage{amsmath,amssymb,mathtools}
\usepackage[scr=rsfs,cal=euler]{mathalfa}
\usepackage{xfrac}
\usepackage[unicode,hidelinks,bookmarks=false]{hyperref}
\newcommand{\ie}{i.e.\ }
\newcommand{\cf}{cf.\ }
\newcommand{\resp}{respectively\ }
\newcommand{\Subsection}{\subsection}
\providecommand{\nobreakdash}{\nobreak\hbox{-}}
\setlength{\emergencystretch}{2em}
\multlinegap0pt
\usepackage{bm}
\usepackage{bbm}
\usepackage{dsfont}
\usepackage{xspace}
\usepackage{cancel}
\usepackage{mathtools}
\usepackage{amsthm}
\makeatletter
\@ifundefined{defi}{\newtheorem{defi}{Defi}}{}
\@ifundefined{lemm}{\newtheorem{lemm}[defi]{Lemma}}{}
\makeatother
\def\inum{{\sqrt{-1}}}
\title[Shannon entropy for harmonic metrics on cyclic Higgs bundles]{Shannon entropy for harmonic metrics on cyclic Higgs bundles}
\author{Natsuo Miyatake}
\date{Source: arXiv:2410.08571v2 (CC BY 4.0)}
\begin{document}
\maketitle

\noindent\emph{Source and licence.} Shannon entropy for harmonic metrics on cyclic Higgs bundles. Version arXiv:2410.08571v2 (CC BY 4.0), distributed under the Creative Commons Attribution 4.0 International licence, as recorded in the arXiv licence metadata for this exact version and shown on the abstract page https://arxiv.org/abs/2410.08571v2. Full rights notice: licenses/arxiv-2410.08571-CC-BY-4.0.txt.\par\smallskip

\noindent\textit{Editorial scope.} Counted unit: the Lemm whose source label is \texttt{MMBB} in the article (source0.tex lines 509--518, source label \texttt{MMBB}), delivered together with the article's own complete original proof of it (source0.tex lines 519--607, one \texttt{proof} environment of 89 lines). No mathematics is written, repaired or re-derived: statement and proof are the article's own text line for line. Required context retained uncounted: Hsigma (lines 486--490); Hsigma2 (lines 486--490); Hsigma3 (lines 486--490); Hs (lines 492--496); Hsigma4 (lines 492--496); Hsigma5 (lines 492--496). Every definition, display and label of those lines is retained unchanged. Omitted from this package and not redistributed: 1--485, 491--491, 497--508, 608--998. Nothing inside a retained excerpt is omitted or altered: every retained body is delivered exactly as the article's lines are written, so no omission receipt of any kind accompanies this package. Citations: the retained excerpts carry no citation command at all, so no bibliography entry of the article is transcribed and no reference is redistributed. Reference adaptation: none was needed and none was made; every reference inside the delivered unit resolves inside it, so no unresolved reference or citation is produced. Printed slips kept literal: no printed word, symbol, hypothesis or number of the counted statement or of its proof was altered. Layout and class: the article's own class is \texttt{article} with packages including amsmath, amssymb, amscd, comment, color, enumerate, cases, mathtools, xy, amsthm, pb-diagram, geometry; the wrapper is the lane's pinned \texttt{amsart} editorial wrapper at 10pt on A4, which restates the theorem-like environments the retained text needs and adds the article's own macro definitions that the retained text needs (\texttt{\textbackslash inum}), with extra wrapper packages (bm, bbm, dsfont, xspace, cancel, mathtools). The delivered numbering is the unit's own. Assets: no retained span contains a figure environment, a graphics-inclusion command, a PDF-page inclusion, a table or a TikZ picture; no image file is copied into this package, the package declares no asset dependency and no image file is redistributed with it. Licence: version arXiv:2410.08571v2 of this article is distributed under the Creative Commons Attribution 4.0 International licence, as recorded in the arXiv licence metadata for this exact version and shown on the abstract page https://arxiv.org/abs/2410.08571v2. A full rights notice is delivered at licenses/arxiv-2410.08571-CC-BY-4.0.txt. Characters: every retained excerpt is pure ASCII, so the delivered engine, XeLaTeX with its default Latin Modern text font, reports no missing character.
\begin{align}
&\inum\Lambda_{H_1}\partial \bar{\partial} \sigma_1\geq 3e^{\sigma_1}-4+e^{-\sigma_2}, \label{Hsigma}\\
&\inum \Lambda_{H_j}\partial\bar{\partial}\sigma_j =3e^{\sigma_j}-3-e^{\sigma_j+\sigma_{j-1}}+e^{-\sigma_{j+1}}\ \text{for $j=2,\dots, n-1$}, \label{Hsigma2} \\
&\inum \Lambda_{H_n}\partial\bar{\partial} \sigma_n=(4-\delta)e^{\sigma_n}-(3-\delta)-e^{\sigma_{n-1}+\sigma_n}. \label{Hsigma3}
\end{align}

\begin{align}
&\inum \Lambda_{H_1}\partial \bar{\partial} \sigma^\prime_1\geq -3+3e^{\sigma_1^\prime}-e^{\sigma_1^\prime+\sigma_2^\prime}, \label{Hs}\\
&\inum \Lambda_{H_j}\partial \bar{\partial} \sigma^\prime_j=-3+3e^{\sigma_j^\prime}+e^{-\sigma_{j-1}^\prime}-e^{\sigma_j^\prime+\sigma_{j+1}^\prime}\ \text{$j=2,\dots, n-2$}, \label{Hsigma4}\\
&\inum \Lambda_{H_{n-1}}\partial \bar{\partial} \sigma^\prime_{n-1}=(3-\delta)e^{\sigma_{n-1}^\prime}-(4-\delta)+e^{-\sigma_{n-2}^\prime}. \label{Hsigma5}
\end{align}

\begin{lemm}\label{MMBB}{\it 
Suppose that $\varphi$ is smooth. Then the following holds:
\begin{align}
&M_j \leq \frac{4j}{2j+1}-\frac{2j-1}{2j+1}M_{j+1}^{-1}\ \text{for $j=1,\dots, n-1$}, \label{M}\\
&M_j B_j\leq B_{j+1}\ \text{for $j=1,\dots, n$}, \label{MB}\\
&M_{n-j}^\prime \leq \frac{4j-(2j-1)\delta}{2j+1-j\delta}-\frac{2j-1-(j-1)\delta}{2j+1-j\delta}M_{n-(j+1)}^{\prime-1} \ \text{for $j=1,\dots, n-2$}, \label{Mp}\\
&M_{n-j}^\prime B_{n-j}^\prime \leq B_{n-(j+1)}^\prime\ \text{for all $j=1,\dots, n-1$}. \label{MpBp}
\end{align}
}
\end{lemm}

\begin{proof}
For each Hermitian metric $H$ on $K_X^{-1}\rightarrow X$, we denote by $\Lambda_{H}$ the contraction operator of the K\"ahler metric induced by $H$. We first prove (\ref{M}) and (\ref{MB}) by induction on $j$. From (\ref{Hsigma}), we have 
\begin{align*}
\inum\Lambda_{H_1}\partial \bar{\partial} \sigma_1\geq 3e^{\sigma_1}-4+M_2^{-1}.
\end{align*}
From the Cheng-Yau maximum principle, we have
\begin{align*}
0\geq 3M_1-4+M_2^{-1}.
\end{align*}
This implies
\begin{align*}
M_1\leq \frac{4}{3}-\frac{1}{3}M_2^{-1} \ \text{and} \ M_1B_1\leq B_2.
\end{align*}
Therefore we obtain (\ref{M}) and (\ref{MB}) for the case where $j=1$.
Suppose that (\ref{M}) holds for some $j=k$, where $1\leq k<n-1$. We show that this implies inequality (\ref{MB}) for $j=k+1$. From (\ref{Hsigma2}), we have 
\begin{align*}
\inum \Lambda_{H_{k+1}}\partial\bar{\partial}\sigma_{k+1} \geq (3-M_k)e^{\sigma_{k+1}}-3+M_{k+2}^{-1}.
\end{align*}
Since we have (\ref{M}) for the case where $j=k$, it holds that $3-M_k>0$. Then from the Cheng-Yau maximum principle, we have
\begin{align*}
0\geq (3-M_k)M_{k+1}-3+M_{k+2}^{-1}.
\end{align*}
This implies 
\begin{align*}
M_{k+1}B_{k+1} \leq B_{k+2}.
\end{align*}
Suppose that (\ref{M}) holds for $j=n-1$. We show that this implies inequality (\ref{MB}) for $j=n$. From (\ref{Hsigma3}), we have
\begin{align*}
\inum \Lambda_{H_n}\partial\bar{\partial} \sigma_n\geq (4-\delta-M_{n-1})e^{\sigma_n}-(3-\delta). 
\end{align*}
Since we have (\ref{M}) for the case where $j=n-1$, it holds that $4-\delta-M_{n-1}>0$. Then from the Cheng-Yau maximum principle, we have 
\begin{align*}
0\geq (4-\delta-M_{n-1})M_n-(3-\delta).
\end{align*}
This implies (\ref{MB}) for the case where $j=n$. Suppose that (\ref{MB}) holds for some $j=k$, where $2\leq k\leq n-1$ and that inequality (\ref{M}) holds for $j=k-1$. We show that this implies (\ref{M}) holds for $j=k$, which completes the proof of (\ref{M}) and (\ref{MB}). From (\ref{MB}) for $j=k$, we have
\begin{align}
M_k(3-M_{k-1})\leq 3-M_{k+1}^{-1}. \label{mk}
\end{align}
By combining (\ref{mk}) with (\ref{M}) for $j=k-1$, we have
\begin{align*}
M_k\left(3-\frac{4k-4}{2k-1}+\frac{2k-3}{2k-1}M_k^{-1}\right)\leq 3-M_{k+1}^{-1}.
\end{align*}
This implies
\begin{align*}
M_k \leq \frac{4k}{2k+1}-\frac{2k-1}{2k+1}M_{k+1}^{-1}.
\end{align*} 
By induction on $j$, we obtain (\ref{M}) and (\ref{MB}) from the above. We next prove (\ref{Mp}) and (\ref{MpBp}). From (\ref{Hsigma5}), we have 
\begin{align*}
\inum \Lambda_{H_{n-1}}\partial \bar{\partial} \sigma^\prime_{n-1}\geq (3-\delta)e^{\sigma_{n-1}^\prime}-(4-\delta)+M_{n-2}^{\prime-1}. 
\end{align*}
Then from the Cheng-Yau maximum principle, we have
\begin{align*}
0\geq (3-\delta)M_{n-1}^\prime-(4-\delta)+M_{n-2}^{\prime-1}.
\end{align*}
This implies
\begin{align*}
M_{n-1}^\prime\leq \frac{4-\delta}{3-\delta}-\frac{M_{n-2}^{\prime-1}}{3-\delta} \ \text{and} \
M_{n-1}^\prime B_{n-1}^\prime\leq B_{n-2}^\prime,
\end{align*}
and thus we have proved (\ref{Mp}) and (\ref{MpBp}) for the case where $j=1$. Suppose that (\ref{Mp}) holds for some $j=k$, where $1\leq k<n-2$. We show that this implies (\ref{MpBp}) for $j=k+1$. From (\ref{Hsigma4}), we have
\begin{align*}
\inum \Lambda_{H_{n-(k+1)}}\partial \bar{\partial} \sigma^\prime_{n-(k+1)}\geq -3+3e^{\sigma_{n-(k+1)}^\prime}+M_{n-(k+2)}^{\prime-1}-M_{n-k}^\prime e^{\sigma_{n-(k+1)}^\prime}
\end{align*}
Since we have (\ref{Mp}) for $j=k$, it holds that $3-M_{n-k}^\prime>0$. Then by the Cheng-Yau maximum principle, we have
\begin{align*}
0\geq -3+3M_{n-(k+1)}^\prime+M_{n-(k+2)}^{\prime-1}-M_{n-k}^\prime M_{n-(k+1)}^\prime.
\end{align*}
This implies $M_{n-(k+1)}^\prime B_{n-(k+1)}^\prime\leq B_{n-(k+2)}^\prime$. Suppose that (\ref{Mp}) holds for $j=n-2$. We show that this implies (\ref{MpBp}) for $j=n-1$. From (\ref{Hs}), we have 
\begin{align*}
\inum \Lambda_{H_1}\partial \bar{\partial} \sigma^\prime_1\geq -3+3e^{\sigma_1^\prime}-e^{\sigma_1^\prime}M_2^\prime.
\end{align*}
Since we have (\ref{Mp}) for $j=n-2$, it holds that $3-M_2^\prime>0$. Then from the Cheng-Yau maximum principle, we have
\begin{align*}
0\geq -3+3M_1^\prime -M_1^\prime M_2^\prime.
\end{align*}
This implies $M_1^\prime B_1^\prime\leq B_0^\prime$. Suppose that (\ref{MpBp}) holds for some $j=k$, where $2\leq k\leq n-2$, and (\ref{Mp}) holds for $j=k-1$. We show that this implies (\ref{Mp}) for $j=k$, which completes the proof of (\ref{Mp}) and (\ref{MpBp}). From (\ref{MpBp}) for $j=k$, we have
\begin{align*}
M_{n-k}^\prime(3-M_{n-(k-1)}^\prime)\leq {3-M_{n-(k+1)}^{\prime-1}}.
\end{align*}
Since we have (\ref{Mp}) for $j=k-1$, it holds that
\begin{align*}
M_{n-k}^\prime\left(3-\frac{4(k-1)-(2k-3)\delta}{2k-1-(k-1)\delta}+\frac{2k-3-(k-2)\delta}{2k-1-(k-1)\delta}M_{n-k}^{\prime-1}\right)\leq {3-M_{n-(k+1)}^{\prime-1}}.
\end{align*}
This implies
\begin{align*}
M_{n-k}^\prime \leq \frac{4k-(2k-1)\delta}{2k+1-k\delta}-\frac{2k-1-(k-1)\delta}{2k+1-k\delta}M_{n-(k+1)}^{\prime-1}.
\end{align*}
Therefore we obtain (\ref{Mp}) for $j=k$. This implies (\ref{Mp}) and (\ref{MpBp}). 
\end{proof}

\end{document}
